% Korean Noether paper 40 — fresh producer translation, unchecked.
% Initial translation authority: ED0005 lines 19142--19789.
% Current scope-rebase authority: ED0011 / NOETH-DE-AUTH-v047-20260808.
% Authority SHA-256: 1A44F967B29972E8F99E5C323A479162AD82A23FC457395915A4BB9DDF51AD41.
% Normalized source span: 83,455 bytes; SHA-256 7965805D3A75C3354C85BC7A3E4725F07BF869A8833FC19D74E32BE369427937.
% Translation only; independent source, Korean-language, render, and certification checks remain external.

\providecommand{\frA}{\mathfrak A}
\providecommand{\frako}{\mathfrak o}
\providecommand{\frB}{\mathfrak B}
\providecommand{\frC}{\mathfrak C}
\providecommand{\frG}{\mathfrak G}
\providecommand{\frL}{\mathfrak L}
\providecommand{\frM}{\mathfrak M}
\providecommand{\frN}{\mathfrak N}
\providecommand{\frR}{\mathfrak R}
\providecommand{\frS}{\mathfrak S}
\providecommand{\frT}{\mathfrak T}
\providecommand{\Gcal}{\mathfrak G}
\providecommand{\mbarA}{\overline{\mathfrak A}}

\section*{40. 비가환 대수}

\begin{center}
\emph{Math. Zs. 37 (1933), 514--541쪽}
\end{center}

\begin{center}
\textbf{비가환 대수학}\\[0.5em]
괴팅겐의 에미 뇌터
\end{center}

가환대수학의 주요 정리들은 갈루아 이론에 들어 있음이 알려져 있다. 그에 앞서는 것은 분리체와 분해체, 곧 주어진 다항식이 일차인자를 하나 분리해 내거나 완전히 일차인자들로 분해되기에 충분한 체들의 이론이다.

여기서는 이에 대응하는 대수학의 부분들을 비가환적으로, 특히 초복소적인 경우에 전개한다. 원리적으로 비가환 방법, 곧 비가환체에서의 표현을 사용한다. 끝에서는 앞에서 말한 가환대수학의 정리들을 가환체에서의 표현을 사용하여 완전히 평행하게 정초할 수 있음을 보인다.

바탕이 되는 표현론 앞에는 짧은 자기동형사상 이론이 놓인다(§ 1). 이 표현론은 표현가군 이론에 바탕을 둔 기존 이론을 발전시킨 것이다\footnote{E. Noether, \emph{Hyperkomplexe Größen und Darstellungstheorie}, Math. Zeitschr. 30 (1929), 641--692쪽. 이하 \emph{표현론}으로 인용한다. van der Waerden, \emph{Moderne Algebra II}의 서술도 보라.}. 특히 직접 표현과 함께 반대 표현을 고찰한다. 한쪽은 다른 쪽으로 환원할 수 있으며, 반대 표현은 이제 반대 표현가군에 의해 생성된다. 장점은 이 반대 표현가군, 곧 쌍가군을 한 확대환 위의 한쪽 가군으로도 볼 수 있다는 데 있다(§ 2). 이로써 비가환체에서의 표현은 확대환의 아이디얼 이론으로 환원된다. 기약 표현류는 확대환의 기약 아이디얼류에 대응하며, 이는 초복소계를 그 가환 계수영역에서 표현할 때 그 계 자체에 성립하는 사실들과 정확히 유사하다(§§ 3, 4).

여기서부터 비가환체 위 행렬환의 구조정리들이 나온다. 핵심은 \emph{모든 부분환은 이 체에 의한 표현, 또는 반대동형인 체에 의한 반대 표현을 뜻한다}는 관찰이다(§ 5). 이 반대동형인 체는 가환인 경우의 최소 분리체·분해체에 대한 첫 유사물이다. 실제로 모든 1차 반대 표현을 매개하고, 중심에 관하여 유한 랭크이면 --- 지금까지는 이를 가정하지 않았다 --- 랭크 1인 직합성분들로의 완전한 분해도 매개한다. 이것이 체의 갈루아 이론을 주며(§ 6), 동시에 반대동형인 나눗셈대수들, 곧 중심에 관하여 유한 랭크인 체들이 R. 브라우어의 대수류군에서 서로 역인 류를 생성한다는 사실을 나타낸다(§ 7).

더 일반적으로 단순계에 성립하며 본문에서는 먼저 제시할 갈루아 이론의 둘째 정초는, 서로 원소별로 교환하는 부분환에 관한 한 정리의 거의 직접적인 귀결이다. 그 정리 자체도 앞서 둔 자기동형사상 고찰에 거의 바로 이어진다.

여기까지 고찰은 순전히 비가환적이다. 그러나 \emph{가환} 분해체의 문제도 § 5에서 먼저 말한 관찰에 따라 분해체를 나눗셈대수 안에 표현함으로써 비가환적으로 다룬다(§ 7). 마지막에는 앞에서 말한 가환적인 경우로의 이행(§ 8)과 임의의 계에 대한 분해체 이론(§ 9)을 제시한다. 이로써 가환체에서의 통상적인 표현론과의 관계도 얻는다.

R. 브라우어는 가환체에서의 이 표현론과 “무리” 인자계에 분해체 이론을 세웠다\footnote{우리의 공동 논문 \emph{Über minimale Zerfällungskörper irreduzibler Darstellungen}, Sitz. Ber. d. Preuß. Ak. d. Wiss. 1927, 221--228쪽을 보라. 222쪽의 주에는 “경위 보고”가 들어 있다. 주요 정리들은 서로 독립적으로 거의 동시에 생겼다. 또 R. Brauer, \emph{Über Systeme hyperkomplexer Zahlen}, § 3, Math. Zeitschr. 30 (1929), 79--107쪽을 보라. A. A. Albert는 뒤에 이 정리들을 독립적으로 다시 발견하였다.}. 이 가환적 정초 때문에 그는 --- 뒤의 Albert도 마찬가지로 --- 중심을 완전체라고 가정해야 했지만, 비가환적 정초에서는 이 제한이 필요 없다. 같은 제한 아래, 역시 가환체에서의 표현론을 사용하여, 나의 비가환체 갈루아 이론을 알게 된 뒤 R. Brauer와 K. Shoda는 이 이론을 더 발전시켰다. R. Brauer는 앞서 말한 교환 부분환 정리를 주었고, K. Shoda는 독립적으로 반단순계까지 포함하는 완전한 이론을 주었다\footnote{R. Brauer, \emph{Über die algebraische Struktur von Schiefkörpern}, § 2, Journ. f. Math. 166 (1932), 241--252쪽. K. Shoda, \emph{Über die galoissche Theorie der halbeinfachen hyperkomplexen Systeme}, Math. Ann. 107 (1932), 252--258쪽. J. Levitzki도 이 결과들을 전개하였다.}.

끝으로 체의 경우에 대해서는 van der Waerden II, § 128에 짧은 서술이 있음을 말해 둔다. 이는 내가 이 문제들을 처음 강의한 1928년 여름학기 강의를 따른 것이다. van der Waerden의 서술은 그 강의보다 여러 곳을 단순화하였고, 나는 그 일부는 두 번째 강의(1929/30년 겨울학기)\footnote{첫 강의는 G. Köthe가, 둘째 강의는 M. Deuring이 정리하였다.}에서, 일부는 비로소 여기에서 받아들여 더 전개하였다. 특히 갈루아 이론의 초복소적 증명법을 가환적인 경우로 옮긴 § 8은 둘째 강의에서 처음 나왔다. 그때 비가환적 증명(§ 6, 3)은 첫 강의보다 본질적으로 단순화되어 있었다. 교환 부분환 정리(§ 5)와 이에 바탕을 둔 비가환 갈루아 이론의 둘째 증명법(§ 6, 1과 2)은 둘째 강의 뒤 R. Brauer의 주 3 논문과 van der Waerden § 128을 알게 된 다음에야 더해졌다. 다만 § 5, 3에서는 거기보다 훨씬 약한 유한성 가정만 쓴다.

첫 강의의 어떤 논법들, 곧 차아이디얼의 교집합을 만드는 방법은 더 복잡하지만 무한 랭크인 계와 정수 계수 계에도 옮길 수 있다는 장점이 있다\footnote{이 방법은 본질적으로 G. Köthe, \emph{Schiefkörper unendlichen Ranges über dem Zentrum}, § 5, Math. Ann. 105 (1931), 15--39쪽에 재현되어 있다. 주 10도 보라. 이 교집합 방법은 이제 양쪽 아이디얼들이 불변가군에 속한다는 거의 자명한 사실(§ 4, 1)로 대체되었다. 이 사실을 처음 알아차린 이는 van der Waerden이다. 이로써 모든 것을 더 단순한 직합분해로 정초할 수 있지만, 무한 랭크인 경우와 정수 계수인 경우에는 이 방법이 성립하지 않는다.}. 가환 정수 계수 계에서는 이로부터 아이디얼 미분과 디퍼런트의 관계에 이른다\footnote{프라하 강연, Jahresber. d. Deutsch. Math. Ver. 39 (1930), 17쪽(비스듬한 쪽매김)을 보라.}. 이에 대해서는 다른 기회에 더 상세히 다루겠다.

분해체 이론의 주된 응용은 교차곱과 그 인자계의 이론에 있으며, 이 이론은 다시 수론적 응용의 기초가 된다. 여기서는 더 다루지 않는다\footnote{교차곱의 이론은 둘째 강의에서 전개하였다. 여기서 준 이론을 전제하지 않도록 조금 고친 꼴이 H. Hasse, \emph{Theory of cyclic algebras over an algebraic numberfield}, 제II장, Transactions of the Amer. Math. Soc. 34 (1932), 171--214쪽에 재현되어 있다. 강의에 더 밀접한 서술은 M. Deuring의 \emph{Ergebnisse der Mathematik} 보고서에 실릴 예정이다.}.

\begin{center}
\textbf{§ 1}\\[0.4em]
\textbf{자기동형사상, 가군, 쌍가군}
\end{center}

표현가군에 바탕을 둔 표현론은 연산자를 갖거나 갖지 않는 아벨군의 자기동형환 이론에 기초한다(§ 2도 보라). 여기서 암묵적으로 쓰는 논법, 곧 사상과 연산법칙 사이의 관계를 반복하지 않도록 몇 가지 간단한 정리로 정식화한다.

\noindent\textbf{1. 곱셈적 사상. 결합법칙.} 먼저 \(\frG\)를 연산자가 없는 군이라 하고, \(\frA\)를 그 절대 자기동형영역, 곧 \(\frG\)에서 자기 자신으로 가는 모든 준동형의 계라 하자. 곱 \(\sigma\tau\)를
\[
g(\sigma\tau)=(g\sigma)\tau
\quad\text{(\(g\in\frG,\ \sigma,\tau\in\frA\))} \tag{1}
\]
로 정의하면 \(\frA\)는 곱셈에 닫혀 있고 결합법칙을 만족한다.

\(\frG\)에 기호 \(\Theta,H,\ldots\)의 집합 \(\frB\)를 주어 연산자 있는 군으로 만든다는 것은, 연산 \(g\Theta,gH,\ldots\)가 \(\frG\)의 일의적으로 정의된 원소를 주고 \(\frG\)의 자기준동형을 생성한다는 뜻이다. 즉 \((g\cdot h)\Theta=g\Theta\cdot h\Theta\)이다. 따라서 \(\frB\)에서 절대 자기동형영역 \(\frA\)의 한 부분집합 \(\overline{\frB}\)로 가는 일의적 사상, 일반적으로는 일대일일 필요가 없는 사상이 존재한다. 앞서 말한 논법은 여기서 다음과 같다.

\emph{연산자영역 \(\frB\)가 곱셈에 닫혀 있으면, \(\frB\)에서 \(\overline{\frB}\)로 가는 일의적 사상이 곱셈적 준동형인 것은 (1)에 대응하는 결합관계}
\[
g(\Theta H)=(g\Theta)H
\quad\text{(\(g\in\frG,\ \Theta,H\in\frB\))} \tag{1a}
\]
\emph{가 성립할 때 그리고 오직 그때이다. 왼쪽 연산자를 다루는 경우에는 이에 대응하여 반대 준동형을 정의한다.}

실제로 \(\frB\)에서 \(\overline{\frB}\)로 가는 사상은 모든 \(g\in\frG\)에 대한 \(g\Theta=g\sigma\)로 주어지므로
\[
(g\Theta)H=(g\sigma)\tau=g(\sigma\tau)
\]
이다. 따라서 (1a)는 곱셈적 준동형의 필요충분조건이다. 왼쪽 연산자가 반대 준동형을 만드는 까닭은 왼쪽에 쓴 자기동형사상 \(\sigma^\ast g,\tau^\ast\sigma^\ast g\)는 오른쪽에서 왼쪽으로 읽어야 하기 때문이다. 곧 \(\tau^\ast\sigma^\ast g=g\sigma\tau\)인 반면, 연산자는 언제나 왼쪽에서 오른쪽으로 읽는다.

\noindent\textbf{2. 연산자준동형 사상.} \(\frG\)가 연산자 있는 군이면 연산자준동형은 잘 알려진 대로
\[
\begin{array}{r@{\quad}c@{\quad}c@{\quad}c@{\qquad}c@{\qquad}r@{\quad}c@{\quad}c@{\quad}c}
(2)&(g\Theta)\sigma&=&(g\sigma)\Theta
&\text{또는}&
(2*)&(\Theta g)\sigma&=&\Theta(g\sigma)
\end{array}
\]
로, 곧 교환적 결합 또는 이어지는 결합법칙으로 정의한다. 두 연산자영역에 1의 논법, 곧 자기동형영역으로의 사상을 적용하면 다음을 얻는다.

\emph{\(\Theta,H,\ldots\)를 원소로 갖는 연산자영역 \(\frB\)와 \(\bar\Theta,\bar H,\ldots\)를 원소로 갖는 연산자영역 \(\frC\)가 주어졌다고 하자. \(\frC\)의 원소들이 \(\frB\)-자기동형사상을 생성하고 동시에 \(\frB\)의 원소들이 \(\frC\)-자기동형사상을 생성하는 것은, 두 영역이 \(\frG\)와 교환적으로 결합되거나 이어지는 결합법칙}
\[
\begin{array}{r@{\quad}c@{\quad}c@{\quad}c@{\qquad}r@{\quad}c@{\quad}c@{\quad}c}
(2a)&(g\Theta)\bar H&=&(g\bar H)\Theta,&
(2a*)&(\Theta g)\bar H&=&\Theta(g\bar H)
\end{array}
\]
\emph{이 성립할 때 그리고 오직 그때이다.}

\noindent\textbf{3. 환 위의 가군과 쌍가군.} 환 \(\frR\) 위의 \emph{오른쪽 가군} \(\frM\)은 \(\frR\)의 원소들을 오른쪽 연산자로 갖는 덧셈 아벨군으로 정의한다. 이때 결합관계 (1a), 연산자를 정의하는 관계 \((g+h)\Theta=g\Theta+h\Theta\)와 함께 분배관계
\[
\begin{array}{r@{\quad}c@{\quad}c@{\quad}c}
(3a)&g(\Theta+H)&=&g\Theta+gH
\end{array}
\]
가 성립해야 한다. 왼쪽 가군도 이에 대응하여 정의한다.

\emph{쌍가군}에는 두 종류가 있다. 두 환 \(\frR,\frS\) 위의 오른쪽 가군을 쌍가군이라 하는 것은, \(\frR,\frS\) 각각에 대한 관계 (1a), (3a)에 더하여 \(\frR,\frS\)가 (2a)에 따라 \(\frM\)과 교환적으로 결합될 때이다. \(\frR\)-왼쪽·\(\frS\)-오른쪽 가군은 다른 관계들에 (2a*), 곧 이어지는 결합법칙이 더해질 때 쌍가군이라 한다.

이 관계들의 의미는 \emph{아벨}군의 경우 절대 자기동형영역이 하나의 환, 곧 절대 자기동형환을 이룬다는 데서 나온다\footnote{비아벨군에서는 “일반화된” 환을 다룬다. 이에 대해서는 H. Fitting의 Math. Annalen 게재 예정 논문을 보라[뒤에 Math. Annalen 107, 514--542쪽에 실림].}. 구체적으로 사상에 관한 위 논법은 다음을 준다.

\emph{덧셈으로 쓴 아벨군 \(\frM\)의 연산자영역 \(\frR\)이 환이면, \(\frR\)에서 절대 자기동형환의 부분집합 \(\overline{\frR}\)으로 가는 일의적 사상이 환 준동형인 것은 \(\frM\)이 \(\frR\) 위 오른쪽 가군, 곧 (1a), (3a)를 만족할 때 그리고 오직 그때이다. 왼쪽 가군이면 반대 환 준동형이다.}

\emph{\(\frM\)이 환 \(\frR,\frS\) 위 오른쪽 가군이면, \(\frM\)이 쌍가군, 곧 (2a)를 만족하는 것은 \(\frR\)을 \(\frS\)-자기동형환의 부분환에, 동시에 \(\frS\)를 \(\frR\)-자기동형환의 부분환에 환 준동형으로 보낼 수 있을 때 그리고 오직 그때이다. \(\frM\)이 \(\frS\)-오른쪽 가군이자 \(\frR\)-왼쪽 가군이면, 쌍가군 조건 (2a*)는 \(\frS\)에서 \(\frR\)-자기동형환의 부분환으로 가는 준동형과 \(\frR\)에서 \(\frS\)-자기동형환의 부분환으로 가는 반대 준동형을 뜻한다.}

마지막 명제는 알려진 귀결들을 다음처럼 정밀화한다.

\emph{1) \(\frR\)이 단위원을 갖는 환이면, \(\frR\)은 \(\frR\)-왼쪽 가군으로서의 자기동형환과 직접 동형이고, \(\frR\)-오른쪽 가군으로서의 자기동형환과 반대동형이다.} 실제로 이어지는 결합법칙 (2a*)가 성립하므로 준동형 명제가 적용된다. 단위원의 존재 때문에 모든 자기동형사상은 \(\frR\)의 어떤 \(a\)에 대한 \(e\mapsto a\) 꼴이다. 따라서 동형이 성립하고 \(\frR\)이 \(\frR\)-자기동형환 전체를 채운다.

\emph{2) \(\frR\)이 일반적으로 비가환인 체 \(A\) 위의 완전행렬환}
\[
\frR=\sum c_{ik}A=\sum A c_{ik}
\]
\emph{이면, \(A\)는 단순 오른쪽 아이디얼들의 자기동형체와 반대동형이고 단순 왼쪽 아이디얼들의 자기동형체와 직접 동형이다.} 실제로 각 단순 오른쪽 아이디얼은 연산자동형인 \(A\)-왼쪽·\(\frR\)-오른쪽 쌍가군이 되고, \(A\)가 \(\frR\)-자기동형사상 전체를 채운다. 왼쪽 아이디얼도 이에 대응한다.

\noindent\textbf{4. 오른쪽 쌍가군에서 곱환 위 한쪽 가군으로 이행.} 오른쪽 쌍가군의 중요성은 본질적으로 이 이행에 바탕을 둔다.

\noindent\textbf{정리.} \emph{\(\frM\)이 두 부분환 \(\frR,\frS\)를 포함하고 이들이 원소별로 서로 교환하는 한 환 \(\frT\) 위 오른쪽 가군이면, \(\frM\)을 \(\frR,\frS\) 위 쌍가군으로도 볼 수 있다. 역으로 \(\frR,\frS\) 위 쌍가군이 주어지고, \(\frR,\frS\)가 원소별로 교환하는 곱환 \(\frT\)가 존재하면, \(\frM\)을 \(\frT\) 위 오른쪽 가군으로도 볼 수 있다. 이때 \(\frT\)에 의한 연산은 주어진 \(\frR,\frS\)에 의한 연산의 연장이다.}

첫 명제는 \(\frT\)가 정의에 따라 절대 자기동형환의 한 부분환 \(\overline{\frT}\)에 준동형으로 사상되고, \(\frR,\frS\)가 원소별로 교환하는 부분환 \(\overline{\frR},\overline{\frS}\)로 옮겨간다는 데서 나온다. 이는 바로 쌍가군을 특징짓는 관계 (2a)이다.

역으로 \(\frM\)이 \(\frR,\frS\)-쌍가군, 더 정확히는 양쪽 모두에 대한 오른쪽 가군으로 주어지면, \(\frR,\frS\)는 절대 자기동형환의 원소별로 교환하는 부분환 \(\overline{\frR},\overline{\frS}\)에 준동형으로 사상된다. 절대 자기동형환 안에는 \(\overline{\frR},\overline{\frS}\)의 곱환 \(\overline{\frT}\)가 존재한다\footnote{여기서 곱환은 공통 상환을 갖는 두 환이 생성하는 환일 뿐이며, 곱이 직합일 필요는 없다. 주어진 교집합에 대해 곱환이 언제나 존재하는 것은 아님을 다음 예가 보인다. \(\frako\)를 정수환이라 하고 \(\frR=\frako[x]\), \(\frS=\frako[y]\), \(2x-1=0\), \(2y-1=0\)이라 하자. \(x,y\) 사이에 연산관계를 정하지 않으면 \(\frako\)가 \(\frR,\frS\)의 교집합이다. 그러나 \(\frR,\frS\)가 한 공통 상환에 들어 있으면 \((2x-1)y-x(2y-1)=0\), 따라서 \(x=y\)이다. 그러므로 교집합이 \(\frako\)인 곱환은 없다.}. 원소별 교환성 때문에 이 환은
\[
\sum \bar r_i\bar s_i+\bar r+\bar s
\]
꼴의 모든 원소로 이루어진다. \(\overline{\frR},\overline{\frS}\)가 단위원을 가지면 덧붙인 항 \(\bar r,\bar s\)는 필요 없다. 이제 \(\frR,\frS\)가 원소별로 교환하는 곱환 \(\frT\)도 존재하면, 이에 따라 \(\frT\)는 \(\sum r_i s_i+r+s\) 꼴의 모든 원소로 이루어진다. 준동형
\[
\frR\to\overline{\frR},\qquad
\frS\to\overline{\frS}
\]
는 대응
\[
\sum r_i s_i+r+s
\longmapsto
\sum\bar r_i\bar s_i+\bar r+\bar s
\]
에 의해 준동형 \(\frT\to\overline{\frT}\)로 연장된다. 따라서 연산
\[
m\Bigl(\sum r_i s_i+r+s\Bigr)
=\sum(mr_i)s_i+mr+ms
\]
는 일의적이며 § 1, 3에 따라 주어진 \(\frR,\frS\)-가군을 포함하는 \(\frT\)-가군을 정의한다.

\begin{center}
\textbf{§ 2}\\[0.4em]
\textbf{반대 표현과 직접 표현}
\end{center}

\noindent\textbf{1. 선형형식 가군.} § 1에서 전개한 이론에서 표현론으로 넘어가는 과정은 거기의 가군들을 선형형식 가군으로 특수화하고 그 자기동형환을 고찰하는 데 바탕을 둔다.

\(\frS\)가 단위원을 갖는 환일 때 \(\frS\)-오른쪽 가군 \(\frM\)을 \(\frS\) 안의 선형형식 가군이라 하는 것은, \(\frM\)이 \(n\)개의 한 항짜리 \(\frS\)-가군의 직합
\[
\frM=m_1\frS+\cdots+m_n\frS
\]
이고 각 \(m_i\frS\)가 \(\frS\)와 연산자동형, 곧 \(m_i s=0\)이면 언제나 \(s=0\)인 경우이다. 이때 \(\frS\)는 절대 자기동형환의 한 부분환에 단지 준동형일 뿐 아니라 동형으로 사상된다.

\noindent\textbf{정리 1.} \emph{\(\frM\)이 \(\frS\) 안의 선형형식 가군이면, \(\frM\)의 \(\frS\)-자기동형환 \(\mA\)는 \(\frS\)의 원소를 성분으로 하는 모든 \(n\)행 행렬의 환 \(\mbarA\)와 단지 준동형일 뿐 아니라 반대동형이다.}

\(\frM\)의 임의의 \(\frS\)-자기동형사상 \(\alpha\)는 대응 \(m_i\mapsto\overline m_i=m_i\alpha\)로 완전히 기술된다. 정의에 따라 이로부터
\[
\sum m_i s_i\longmapsto
\sum\overline m_i s_i
=\sum(m_i\alpha)s_i
=\sum(m_i s_i)\alpha
\]
가 나오기 때문이다. 행렬표기로
\[
(m_1\alpha,\ldots,m_n\alpha)=(m_1,\ldots,m_n)A
\]
라 쓰면, \(\alpha\)가 모든 자기동형사상을 두루 돌 때 대응 \(\alpha\mapsto A\)는 \(\mA\)와 \(\mbarA\) 사이의 일대일 대응을 준다. 실제로 \(\frM\)을 선형형식 가군이라고 가정했으므로 \(A\)는 \(\alpha\)로 유일하게 정해지고, \(\frS\)의 원소를 성분으로 하는 모든 \(n\)차 행렬 \(A\)는 자기동형사상을 생성한다. 또한
\[
\alpha+\beta\mapsto A+B,\qquad
\alpha\beta\mapsto BA
\]
이다. 마지막 것은
\[
(m_1\alpha\beta,\ldots,m_n\alpha\beta)
=(m_1,\ldots,m_n)A\beta
=(m_1,\ldots,m_n)\beta A
\quad[\text{\(ms\beta=m\beta s\)이므로}]
=(m_1,\ldots,m_n)BA
\]
에서 나온다.

\noindent\textbf{2. 표현과 표현가군.} 환 \(\frR\)의 \(\frS\) 안 \(n\)차 반대 표현 또는 직접 표현을, \(\frR\)에서 \(\frS\)의 원소를 성분으로 하는 모든 \(n\)행 행렬의 환의 한 부분환으로 가는 반대 환 준동형 또는 직접 환 준동형으로 이해하면, 정리 1은 다음과 같이도 쓸 수 있다.

\noindent\textbf{정리 1'.} \emph{\(\frS\) 안 \(n\)차 선형형식 가군의 자기동형환 \(\mA\)는 \(\frS\)의 모든 \(n\)행 행렬로 이루어진 완전행렬환 \(\mbarA\)에 의한 동형인, 곧 충실한 반대 표현을 갖는다.}

여기서부터 직접 표현에 관한 통상적인 정리들이 반대 표현에도 성립한다.

\noindent\textbf{정의.} \(\frS\) 안의 선형형식 가군 \(\frM\)이 \(\frR,\frS\) 위 쌍가군이면서 동시에 \(\frR\)-오른쪽 가군이자 \(\frS\)-오른쪽 가군이면 \(\frR\)의 \(\frS\) 안 반대 표현가군이라 한다. \(\frM\)이 쌍가군이면서 \(\frR\)-왼쪽·\(\frS\)-오른쪽 가군이면 직접 표현가군이라 한다.

\noindent\textbf{정리 2.} \emph{각 반대 표현가군 또는 직접 표현가군은 \(\frR\)의 \(\frS\) 안 동치인 반대 표현 또는 직접 표현들의 한 류를 생성하며, 모든 반대 표현과 직접 표현은 이렇게 생성된다.}

먼저 \(\frM\)이 반대 표현가군이라 하자. § 1, 3에 따라 \(\frR\)은 \(\frM\)의 \(\frS\)-자기동형환의 한 부분환 \(\overline{\frR}\)에 직접 준동형으로 사상된다. § 2, 정리 1'에 따라 \(\overline{\frR}\)은 충실한 반대 표현을 가지며, 이는 \(\frR\)의 반대 표현이기도 하다. 역으로 반대 표현 \(\frR\to\frR^*\)가 주어지면, 이를 \(\frR^*\)에서 \(\frS\)-자기동형환의 부분환 \(\overline{\frR}\)으로 가는 반대동형과 합성하여 \(\frR\)에서 \(\overline{\frR}\)으로 가는 직접 준동형을 얻는다. 따라서 \(\frM\)은 \(\frR\)-가군이 되고, § 1, 3에 따라 쌍가군, 곧 표현가군이 된다. 다른 \(\frS\)-기저로 바꾸면 동치인 표현들의 류가 나온다.

직접 표현가군의 경우에는 \(\frR\)이 \(\overline{\frR}\)에 반대 준동형으로 사상된다. 이를 \(\overline{\frR}\)의 반대 표현과 합성하면 \(\frR\)의 직접 표현을 얻는다. 역으로 직접 표현에서 \(\overline{\frR}\)으로 가는 반대 준동형 사상이 나오며, § 1, 3에 따라 \(\frM\)은 직접 표현가군이 된다.

\noindent\textbf{주의.} 물론 두 종류의 표현은 서로 환원할 수 있다. \(\frR\)의 직접 표현은 \(\frR\)의 반대환의 반대 표현이고 그 역도 성립한다. 다른 환원은 \(\frS\)에서 반대환으로 옮기고 행렬들을 전치행렬로 바꾸는 것이다. 이는 \(\frR\)-왼쪽·\(\frS\)-오른쪽 가군에서 \(\frR,\frS\)-왼쪽 가군으로 옮기는 데 대응한다.

\noindent\textbf{3. 반대 표현가군에서 확대환 위 가군으로 이행.} 반대 표현가군의 장점은 무엇보다 § 1, 4에 따라 확대환으로 이행할 수 있다는 데 있다. 지금의 특수한 경우, 곧 \(\frM\)이 \(\frS\) 위 선형형식 가군인 경우에는 표현가군에 대한 다음 이행정리가 된다.

\emph{\(\frM\)이 두 부분환 \(\frR,\frS\)를 포함하고 이들이 원소별로 서로 교환하는 환 \(\frT\) 위 오른쪽 가군이며, \(\frM\)이 \(\frS\) 위 선형형식 가군이면, \(\frM\)을 \(\frR\)의 \(\frS\) 안 반대 표현가군으로도 볼 수 있다.}

\emph{역으로 \(\frR\)의 \(\frS\) 안 반대 표현가군이 주어지고, \(\frR,\frS\)가 원소별로 교환하는 곱환 \(\frT\)가 존재하면, \(\frM\)을 \(\frT\)-가군으로도 볼 수 있다. 이때 \(\frT\)-연산은 주어진 \(\frR,\frS\)-연산의 연장이다.}

표현론에서 본질적으로 쓰는 것은 다음 귀결이다.

\noindent\textbf{표현가군 이행정리의 귀결.}\\*
\(\frM\)이 \(\frR\)의 \(\frS\) 안 반대 표현가군이면서 곱환 \(\frT\) 위 가군이기도 하면, \(\frM\)에서 한 가군 \(\frN\)으로 가는 \(\frR,\frS\)-동형은 \(\frT\)-동형과 같다. \emph{따라서 동형인 \(\frT\)-가군들의 각 류에는 \(\frR\)의 \(\frS\) 안 반대 표현류 하나가 대응하고 그 역도 성립한다.}

이 정리들을 § 1, 2와 3에 결합하면 한 표현과 교환하는 행렬들과 \(\frR,\frS\)-자기동형사상들 사이의, 이하에서 기초가 되는 관계를 얻는다.

\noindent\textbf{교환행렬에 관한 정리.}\\*
\emph{\(\frR\to\frR^*\)가 \(\frR\)의 \(\frS\) 안 \(n\)차 반대 표현이고, \(\frB^*\)가 \(\frR^*\)의 모든 원소와 교환하는 \(\frS\) 안 \(n\)차 행렬들의 환이면, \(\frB^*\)는 이 표현을 생성하는 표현가군의 \(\frR,\frS\)-자기동형사상들, 곧 동시에 \(\frR\)-자기동형사상인 \(\frS\)-자기동형사상들의 반대동형 표현을 준다. 곱환 \(\frT\)가 존재하면 이는 \(\frT\)-자기동형사상들의 표현이기도 하다.}

\(\frB\)를 \(\frM\)의 모든 \(\frR,\frS\)-자기동형사상들의 환이라 하자. \(\frS\)-자기동형환 \(\mA\)의 부분환으로서 \(\frB\)는 \(\frS\) 안 반대동형 표현을 갖는다. § 1, 2와 3에 따라 \(\frB\)는 \(\mA\)의 자기동형사상 가운데 \(\frR\)과 교환적으로 결합된 것, 곧 관계 (2a)를 만족하는 것들로 이루어진다. 따라서 \(\frB\)는 \(\overline{\frR}\)의 각 원소와 교환하는 자기동형사상 전체이며, 여기서 \(\overline{\frR}\)은 \(\mA\) 안의 \(\frR\)의 상이다. \(\frB,\overline{\frR}\)의 표현은 단순한 반대 준동형이 아니라 반대동형이므로 이것이 증명할 사실이다.

\noindent\textbf{주의.} \(\frR,\frS\)-자기동형사상은 대응 \(m_i\mapsto m_i'\)를 뜻하며, \(m_i'\)들에 의해 같은 표현 \(\frR^*\)이 생긴다. 다만 \(\frB^*\)의 행렬들이 반드시 가역은 아니라는 사실에 대응하여, \(m_i'\)들이 반드시 \(\frM\)의 \(\frS\)-기저를 이루지는 않는다. 합 \(\sum m_i'\frS\)가 더 이상 직합이 아닐 때 “표현”은 확장된 뜻으로 이해해야 한다.
\[
(m_1',m_2',\ldots,m_n')a=(m_1',\ldots,m_n')A
\]
는 여전히 맞지만 \(A\)는 더 이상 \(a\)로 유일하게 정해지지 않는다.

이행정리에 관해 한 가지를 더 지적한다. 대각행렬 \(E\cdot s\)는 \(\frS\)의 \emph{직접 동형} 표현, 곧 \(\frS\)의 “항등” 표현을 준다. 이것이 직접 동형인 까닭은 \(\frM\)의 \(\frS\)-자기동형환이 \(\frS\)와 반대동형인 부분환을 가지며, 지금의 표현은 그 부분환의 반대 표현이기 때문이다. 실제로 \(\frM\)의 각 한 항짜리 성분 \(m_i\frS\)는 오른쪽 가군으로서 \(\frS\)와 연산자동형이고, 따라서 그 \(\frS\)-자기동형환은 \(\frS\)와 반대동형이다(§ 1, 3, 귀결 1).

\(\frR,\frS\)로 일의적으로 정의되는, \(\frT\)에서 \(\frS\) 안의 행렬들로 가는 대응은 \emph{환 \(\frT\)의 표현이 아니라, \(\frR\)의 반대 표현을 \(\frS\)에 관하여 연산자준동형인 사상으로 확대한 것}이다.
\[
\sum r_i s_i\longmapsto\sum R_i s_i.
\]
특히 \(\frR\)의 반대 표현은 동시에 \emph{\(\frR\)의 중심에 놓이는 교집합 \([\frR,\frS]\)에 관하여 연산자준동형}이다.

\begin{center}
\textbf{§ 3}\\[0.4em]
\textbf{한 체에 관한 가군}
\end{center}

이하에서는 일반적으로 비가환인 체 안의 표현만 다룰 것이므로, 한 체 위 선형형식 가군에 관한 몇 가지 간단한 정리를 모아 둔다.

\noindent\textbf{1. 전체 가군의 주어진 기저에 관한 부분가군의 정규기저.} \(A\)를 체, \(\frN=x_1A+\cdots+x_nA\)를 랭크 \(n\)인 선형형식 가군이라 하고, \(\frL=z_1A+\cdots+z_\ell A\)를 랭크 \(\ell\le n\)인 부분가군이라 하자. \emph{\(z_i\)들이 다음 꼴이면 \(x\)에 관한 정규기저라 한다.}
\[
z_i=x_i-(x_{\ell+1}\alpha_{i,\ell+1}+\cdots+x_n\alpha_{i,n}),
\qquad(i=1,\ldots,\ell).
\]
\emph{모든 가군 \(\frL\)은 \(x\)들의 번호를 알맞게 바꾸면 정규기저를 갖는다.}

실제로 번호를 알맞게 바꾸면
\[
\frN=\frL+x_{\ell+1}A+\cdots+x_nA
\]
이고, 따라서
\[
x_i\equiv x_{\ell+1}\alpha_{i,\ell+1}+\cdots+x_n\alpha_{i,n}
\quad(\frL),\qquad(i=1,\ldots,\ell). \tag{2}
\]
그러므로 \(\ell\)개 원소 \(z_i=x_i-(x_{\ell+1}\alpha_{i,\ell+1}+\cdots)\)는 \(\frL\)에 속한다. 이들은 \(x_{\ell+1},\ldots,x_n\)과 함께 \(\frN\)의 기저를 이루므로 \(\frL\) 전체를 생성한다.

\noindent\textbf{2. 확대가군.} 다시 \(A\)를 체, \(P\)를 부분체라 하자. 랭크 \(n\)인 \(A\)-가군 \(\frN\)이 한 \(P\)-가군 \(\frM\)의 \emph{확대가군}, \(\frN=\frM_A\), 이라는 것은 \(\frM\)이 \(\frN\)의 \emph{같은 랭크}를 갖는 부분가군인 경우이다. \(\frM=y_1P+\cdots+y_nP\)이면
\[
\frN=\frM_A=y_1A+\cdots+y_nA.
\]
역으로 각 \(P\)-가군 \(\frM=y_1P+\cdots+y_nP\)에는 가군 동형까지 유일한 확대가군 \(\frM_A\)가 존재한다. 실제로 가군 \(\bar y_1A+\cdots+\bar y_nA\)는 \(\frM\)과 동형인 부분가군 \(\overline{\frM}\)을 포함하며, 이를 \(\frM\)과 동일시할 수 있다.

\noindent\textbf{귀결 a).} \(z_1,\ldots,z_t\)가 \(\frM\)의 \(P\)에 관하여 선형독립인 원소들이면, \(\frN=\frM_A\)에서도 \(A\)에 관하여 선형독립이다. 실제로 이들을 \(\frM\), 따라서 \(\frN\)의 기저로 보충할 수 있다. 그러므로 \(\frM\)의 각 \(P\)-가군 \(\frT=z_1P+\cdots+z_tP\)는 \(\frM_A\) 안에서 확대가군 \(\frT_A=z_1A+\cdots+z_tA\)를 생성한다.

\noindent\textbf{귀결 b).} \(\frM\)의 각 \(P\)-가군 \(\frT=z_1P+\cdots+z_tP\)는 수축가군, 곧 그 확대가군 \(\frT_A\)와 \(\frM\)의 교집합이다.
\[
\frT=\frT_A\cap\frM.
\]
실제로 \(\frT\subseteq\frT_A\cap\frM\)이다. 역으로 \(a\in\frT_A\cap\frM\)이면 \(a=\sum z_i\alpha_i\)이므로 \(\frM\)의 원소 \(a,z_1,\ldots,z_t\)는 \(A\)에 관하여 종속이고, 귀결 a에 따라 \(P\)에 관해서도 종속이다. 따라서 \(a\in\frT\)이다.

\noindent\textbf{3. 불변가군에 관한 정리.} 다시 \(\frN=\frM_A\)를 랭크 \(n\)인 \(P\)-가군 \(\frM\)의 확대가군이라 하자. 또 \(P\)가 \(A\)의 환 자기동형사상들로 이루어진 한 군 \(\frG\)에 대한 완전 불변체라 하자. \(\frG\)는 다음 규정으로 \(\frN=\frM_A\)의 연산자영역이 된다.
\[
G\cdot m=m\quad(m\in\frM,\ G\in\frG),
\]
\[
G\cdot\sum m_i\alpha_i=\sum m_i\,G(\alpha_i)
\quad(m_i\in\frM,\ \alpha_i\in A).
\]

\noindent\textbf{보조정리.} \emph{이 규정 아래 \(\frM\)은 \(\frN\)의 원소 가운데 \(\frG\)에 의해 고정되는 모든 원소로 이루어진다.} 실제로 \(x_1,\ldots,x_n\)이 \(\frM\)의 \(P\)-기저이고 \(w=x_1\alpha_1+\cdots+x_n\alpha_n\)이면, 모든 \(G\in\frG\)에 대한 \(G(w)=w\)에서 \(G(\alpha_i)=\alpha_i\)가 나오고, 가정에 따라 \(\alpha_i\in P\)이다.

\noindent\textbf{정리.} \emph{\(\frG\)를 연산자영역으로 볼 때 허용되는 부분가군은 \(\frM\)의 부분가군 \(\frL\)들의 확대가군 \(\frL_A\)와 오직 그것들이다.}

한쪽 방향은 분명하다. 역으로 \(\frL\)이 \(\frG\)에 관하여 허용된다고 하자. \(\frL\)의 \(x\)에 관한 정규기저를 \(z_1,\ldots,z_\ell\)이라 하자(1을 보라). 여기서 \(x\)는 \(\frM\)의 \(P\)-기저로 택하여 \(\frG\)가 원소별로 고정한다. 가정에 따라
\[
G(z_i)=x_i-\bigl(x_{\ell+1}G(\alpha_{i,\ell+1})
              +\cdots+x_nG(\alpha_{i,n})\bigr)
\]
는 모든 \(G\in\frG\)에 대해 \(\frL\)의 원소이다. 따라서 고정한 \(G\)에 대하여 \(z\)들의 선형결합으로 나타난다. \(x_1,\ldots,x_\ell\)의 계수를 비교하면 모든 \(G\in\frG\)에 대해 \(G(z_i)=z_i\)이다. 보조정리에 따라 \(z_i\in\frM\)이고, \(\frL\)은 \(\frM\)의 \(P\)-가군 \(\frT=z_1P+\cdots+z_\ell P\)의 \emph{확대}이다. 2에 따라 \(\frT=\frL\cap\frM\)이다\footnote{이 절의 정리들은 간단한 정렬정리 논법에 의해 \(\frM\)의 \(P\)-랭크가 무한인 경우에도 유지된다. G. Köthe, \emph{Ein Beitrag zur Theorie der kommutativen Ringe ohne Endlichkeitsvoraussetzung}, § 1, Gött. Nachr. 1931, 195--207쪽을 보라.}.

\begin{center}
\textbf{§ 4}\\[0.5em]
\textbf{초복소계와 그 표현류}
\end{center}

\noindent\textbf{1. 확대정리.} 이제 § 2의 표현정리들을 § 3의 가군정리들과 결합한다. 일반 환 \(\frR\) 대신 가환체 \(P\) 위에서 단위원을 갖는 \emph{초복소계} \(S,T,\ldots\), 곧
\[
S=x_1P+\cdots+x_nP
\]
를 고찰한다. 임의의 표현환 \(\frS\) 대신에는 체 \(A,B,\ldots\)만 고찰하며, 달리 말하지 않으면 그 중심은 \(P\)라고 한다. 이 경우 \(S,A\)가 원소별로 교환하고 교집합이 \(P\)인 곱환이 언제나 존재한다. 이는 직접곱 \(S\times A\)이며 동시에 § 3의 뜻에서 \(S\)의 확대가군
\[
S_A=x_1A+\cdots+x_nA
\]
이므로 확대환이라고도 한다.

이 특수화 아래 불변가군에 관한 정리(§ 3, 3)는 다음 정리가 된다.

\noindent\textbf{확대정리.} \emph{\(S_A\)의 각 양쪽 아이디얼은 \(S\)의 한 아이디얼의 확대아이디얼이며, 더 정확히는 그 아이디얼과 \(S\)의 교집합을 확대한 것이다.}

실제로 바탕에 둔 군 \(\frG\)가 \(A\)의 내부자기동형사상군이면, \(A\)의 중심인 \(P\)는 완전 불변체이다. \(A\)와 \(S\)의 원소별 교환성은 § 3, 3에 따라 \(\frG\)가 \(S_A\)로 확장되고 \(S\)가 완전 불변영역이 된다는 뜻이다. \(S_A\)의 각 양쪽 아이디얼 \(\mathfrak a\)는 \(\alpha^{-1}\mathfrak a\alpha\le\mathfrak a\)이므로 허용되는 부분군이다. 따라서 § 3, 3에 따라 교집합가군 \(\mathfrak a\cap S\)의 확대이며, 이 교집합은 \(S\)의 아이디얼이다.

\noindent\textbf{보충.} \(S\)가 가환이면 \(S\)는 \(S_A\)의 중심이 된다. 실제로 \(S\)는 중심에 있고 \(A\)가 생성하는 내부자기동형사상들에 대한 완전 불변영역이다. \emph{더 일반적으로 \(S\)의 중심은 \(S_A\)의 중심이기도 하다.}

\noindent\textbf{표기.} \(A_r,B_n,\ldots\)는 언제나 각각 \(A,B,\ldots\) 위의 표시한 차수의 행렬환을 뜻한다. 곧
\[
A_r=\sum A c_{ik}=\sum c_{ik}A.
\]

이하에서는 단순계에 대한 확대정리의 귀결을 본질적으로 쓴다. \(S\)가 단순계이면\footnote{단순계는 통상대로 “양쪽 단순”을 뜻한다.} \(S_A\)도 단순이다.
\[
S_A=\sum Bc_{ik}=\sum c_{ik}B=B_t,
\]
여기서 대응하는 체 \(B\)의 중심은 \(P\)를 포함한다. \(A\), 따라서 \(B\)도 \(P\) 위 무한 랭크일 수 있으며, \(S\)만 초복소적이라고 가정한다. \(S,A\)의 공통 중심이 \(P\)이면 \(P\)는 \(S_A\)의 중심이기도 하다.

더 일반적으로 \(S\)가 단순 초복소계이면 직접곱 \(S\times A_r\)도 단순이다. 실제로 \(S\times A_r\)는 \(S_A\times P_r\)와 일치하므로 \(S_A=B_t\)이면 \(B_{tr}\)와 같다.

\noindent\textbf{2. 표현류.} 초복소계의 반대 표현 또는 직접 표현은 통상대로 계수영역 \(P\)에 관하여 연산자준동형이고(§ 2의 끝) 영표현과 다른 것을 뜻한다. 이행정리 및 그 귀결(§ 2, 3)과 1의 확대정리를 결합하면 다음을 얻는다.

\noindent\textbf{표현류에 관한 정리.} \emph{\(S\)가 단위원을 갖고 \(P\)를 계수영역으로 하는 초복소계이며 \(A\)가 중심 \(P\)인 체이면, \(S\)가 갖는 \(A\) 안의 서로 다른 기약 반대 표현류의 수는 \(S_A\)를 그 라디칼로 나눈 잉여류환에서 단순 오른쪽 아이디얼들의 류 수와 같다. \(S\)가 단순계이면 \(A\) 안에 기약 반대 표현류 하나와 기약 직접 표현류 하나를 정확히 갖는다.}\footnote{\normalfont \(A\)가 \(S\)에 딸린 자기동형체인 경우에는 \emph{표현론}, 주 20을 보라.}

실제로 이행정리의 귀결(§ 2, 3)에 따라 단순 반대 표현류와 \(S_A\) 위 단순 가군류는 일대일로 대응한다. \(S_A\)는 체 \(A\)에 관하여 유한 랭크이므로 한쪽 아이디얼에 대한 최대·최소조건을 만족하며, 따라서 명제의 첫 부분은 \emph{표현론} § 19에서 바로 나온다. \(S\)가 단순이면 확대정리에 따라 \(S_A\)도 단순이다. 최대·최소조건 때문에 동시에 완전가약이고, 따라서 단순 한쪽 아이디얼류가 하나뿐이다. 곧 \(S\)는 \(A\) 안에 기약 반대 표현류가 하나뿐이다. \(S\)와 반대동형인 계도 단순이고 기약 반대 표현류가 하나뿐이므로, \(S\)에는 단순 직접 표현류도 하나뿐이다.

\noindent\textbf{3. 랭크 관계.} \(S\)가 단순일 때 \(S_A=\sum Bc_{ik}\)가 \(A\)와 \(B\) 양쪽 모두에 관하여 유한 랭크라는 사실은 다음 랭크 관계를 준다.
\[
\begin{array}{rl}
(1)&n=rt,\quad
n=(S:P),\quad t^2=(S_A:B),\quad
r=\text{기약 반대 표현의 차수}.
\end{array}
\]

실제로 \(S_A\) 자체가 \(S\)의 \(A\) 안 반대 표현가군이다. 그 표현은 이미 \(P\)에 들어 있다. \(B\) 위 행렬환으로서 \(S_A\)는 \(t\)개의 단순 오른쪽 아이디얼, 곧 \(S_A\)-가군으로 분해되며, 각 아이디얼의 \(A\)-랭크는 기약 반대 표현의 차수 \(r\)과 같다. \(n\)은 동시에 \(S_A\)의 \(A\)-랭크이므로 랭크 관계가 따라온다.

\paragraph{4. \(A\)가 \(P\)에 관하여 유한 랭크일 때 랭크 관계의 강화.}
이 경우 세 관계
\[
\begin{array}{rl}
(1)&(S:P)=n=rt,\\
(2)&(B:P)\cdot t=(A:P)\cdot r,\\
(3)&(S:P)\cdot(B:P)=(A:P)\cdot r^2
\end{array}
\]
가 성립한다. 여기서 (2)는 \(S_A\)의 단순 오른쪽 아이디얼의 \(P\)-랭크를 3에 따라 \(B\)-랭크와 \(A\)-랭크로 각각 나타낸 것이다. (3)은 행렬환 \(B_n\) 안의 단순 오른쪽 아이디얼의 랭크를 나타내며, (1)을 사용하여 (2)에 \(r\)을 곱하면 나온다.

\begin{center}
§ 5

\textbf{행렬환에 표현정리 적용}
\end{center}

§ 4에서 고찰한 \(S\)의 \(A\) 안 반대 표현은 \(S\)에서 \(\overline A_r\)의 한 부분환으로 가는 직접 준동형으로도 볼 수 있다. 여기서 \(\overline A\)는 \(A\)와 반대동형인 체이다. \emph{이하 여러 절의 원리는 역으로 행렬환 \(\overline A_r\)의 연구를 반대동형인 체 \(A\) 안의 표현론으로 환원하는 것이다.}

\noindent\textbf{1. \(\overline A_r\) 안의 가약·기약 매장.} \(A,\overline A\)를 중심 \(P\) 위에서 유한 또는 무한 랭크인 반대동형체라 하자\footnote{일반적으로 반대인 체 또는 계는 서로 대응하는 그리스 문자와 라틴 문자로 나타낸다.}. \emph{\(P\) 위 단순계 \(S\)가 \(\overline A_f\) 안에 기약 또는 가약으로 매장 가능하다는 것은, \(S\)가 \(A\) 안에 \(f\)차 기약 또는 가약 반대 표현을 갖는다는 뜻이다.} \(S\)가 \(\overline A_r\) 안에 기약으로 매장 가능하면, \(s>1\)인 모든 행렬환 \(\overline A_{rs}\) 안에는 가약으로 매장 가능하다. 실제로 이 차수들과 오직 이 차수들의 가약 표현을 갖는다(§ 4, 2). 따라서 \(S\)는 \(\overline A_r,\overline A_{rs}\)의 어떤 부분환들과 직접 동형이다. 이 동형은 \(P\)의 항등사상의 연장이다\footnote{표현의 경우와 마찬가지로 이하의 동형은 언제나 \(P\)의 항등사상의 연장, 곧 \(P\) 위의 동형이다. 매번 따로 밝히지는 않는다.}. § 4, 3에 따라 \(r\)은 \(S\)의 \(P\)-랭크 \(n\)의 약수이다.

이 방법으로 \(P\)를 포함하고 \(P\) 위 유한 랭크인, 여러 \(\overline A_f\), \(f=1,2,\ldots\), 의 모든 단순 부분환을 얻는다. 실제로 그러한 각 부분환은 \(A_f\)의 한 부분환과 반대동형이므로 \(A\) 안에 가약 또는 기약 반대 표현을 갖는다.

\noindent\textbf{2. 내부자기동형사상에 관한 정리.} \emph{\(S^{(1)},S^{(2)}\)가 \(\overline A_f\)의 두 단순 부분환으로 \(P\)를 포함하고 \(P\) 위 유한 랭크이며, \(S^{(1)},S^{(2)}\)가 \(P\) 위에서 동형이면, 이 동형은 \(\overline A_f\)의 한 내부자기동형사상으로 생성된다.}

실제로 \(S^{(1)},S^{(2)}\)는 한 단순계 \(S\)를 \(\overline A_f\)에 일반적으로 가약으로 매장한 것, 따라서 \(A\) 안의 \(f\)차 직접 표현들이다. 이들은 \(A\) 안의 같은 가약 직접 표현류에 속하므로, 둘을 옮기는 행렬이 자기동형사상을 생성한다.

\(A\) 자체가 \(P\) 위 유한 랭크여서 초복소적이면 이 정리는 다음의 알려진 정리가 된다.

\emph{\(A_f\)에서 중심 \(P\)를 포함하는 두 단순 부분계가 \(P\) 위에서 동형이면, \(A_f\)의 한 내부자기동형사상에 의해 서로 옮겨간다. 특히 \(A_f\)의 모든 자기동형사상은 내부자기동형사상이다.}\footnote{\normalfont \(A\)가 \(P\) 위 무한 랭크이면 이는 더 이상 성립하지 않는다. 주 5에서 인용한 Köthe의 논문을 보라.}

\noindent\textbf{3. 원소별로 교환하는 부분환에 관한 정리.}
이 정리도 § 2, 3의 교환행렬 정리에서, 이 절의 처음에 말한 원리에 따라 나온다.

\emph{\(S\)가 \(\overline A_f\)의 단순 부분계로 \(P\)를 포함하고, \(R\)이 \(\overline A_f\)에서 \(S\)의 모든 원소와 교환하는 원소들의 전체이면, \(R\)도 행렬환이다. 곧 \(R=\overline B_s\)이다. \(S_A\)에 딸린 체 \(B\)와 \(R\)에 딸린 체 \(\overline B\)는 반대동형이다. \(R,S\)의 교집합은 \(S\)의 중심이다. \(R\)이 체인 것은 \(S\)가 \(\overline A_f\) 안에 기약으로 매장되었을 때 그리고 오직 그때이다.}

실제로 § 2, 3에 따라 \(R\)은 \(S\)의 \(A\) 안 반대 표현가군 \(\frM\)의 자기동형환과 직접 동형이다. 이는 \(\frM\)을 \(S_A\)-가군으로 본 자기동형환이다. \(\frM\)이 \(s\)개의 단순 \(S_A\)-가군으로 분해되고, § 4, 1처럼 \(S_A=\sum Bc_{ik}\)라 두면, \(R\)은 \(\sum\overline Bc_{ik}\)와 동형이다. 여기서 \(B,\overline B\)는 반대동형이다. 실제로 \(B\)는 \(S_A\)의 단순 오른쪽 아이디얼, 곧 단순 \(S_A\)-가군의 자기동형체와 반대동형이다(§ 1, 3의 끝). \(R\)이 체 \(R=\overline B\)인 것은 \(s=1\), 곧 \(S\)가 기약으로 매장된 경우와 정확히 같다. \(R,S\)의 교집합이 \(S\)의 중심이라는 사실은 \(R\)의 정의에서 곧바로 나온다.

\noindent\textbf{4. 행렬환이 초복소적인 경우의 강화된 교환정리.} 이 경우 § 4, 4의 랭크 관계에서 다음 강화가 나온다.

\emph{중심 \(P\) 위 유한 랭크인 체 \(A\)에 대하여, \(\overline A_f\)에서 \(P\)를 포함하는 단순 부분계들은 쌍 \(S,\overline S\)로 나뉜다. \(\overline S\)는 \(S\)의 모든 원소와 교환하는 원소 전체이고 그 역도 성립한다. \(S_A\)와 \(\overline S\)에 딸린 체들은 반대동형이며, \(S\)와 \(\overline S_A\)에 딸린 체들도 반대동형이다. \(S,\overline S\)의 교집합은 공통 중심과 같다. 한 부분환이 체인 것은 다른 부분환이 기약으로 매장되었을 때 그리고 오직 그때이다. \(S,\overline S\)의 \(P\)-랭크의 곱은 \(\overline A_f\)의 랭크와 같다. \(f=rs=\bar r\,\bar s\)이고 \(S\)가 \(\overline A_r\)에 기약으로, \(\overline S\)가 \(\overline A_{\bar r}\)에 기약으로 매장 가능하면, \(S,\overline S\)에 딸린 체들은 \(f=g s\bar s\)인 \(A_g\) 안의 서로 교환하는 부분환 한 쌍과 동형이다.}

본질적으로는 랭크들의 곱에 관한 명제만 증명하면 된다. 먼저 \(S\)가 기약으로 매장되어 \(\overline S\)가 체이고 \(B\)와 반대동형이라고 하자. 여기서 \(S_A=B_t\)이다. 이제 \(f=r\)이고 \(\overline S\)의 랭크는 \(B\)의 랭크와 같으므로 § 4, 4의 랭크 관계 (3)이 곱에 관한 명제를 준다. \(S\)가 가약으로 매장되었으면 \(\overline S=\overline B_s\), \(f=rs\)이고, (3)에 \(s^2\)을 곱하면 명제를 얻는다. 이로부터 \(S\)가 \(\overline S\)의 모든 원소와 교환하는 원소 전체임이 곧바로 나오고, 마지막 명제를 제외한 나머지는 3에서 모두 따라온다.

마지막 명제는 기약 매장으로 두 번 이행하여 얻는다. \(f=rs\)인 \(\overline A_r\) 안에서는 \(S\)가 기약으로 매장되고 \(S,R\)은 서로 교환하는 계이며, \(R\)은 체 \(B\)와 동형이다. 한편 \(S,\overline S\)의 반대성 때문에 \(S=\overline B_s\)이다. 여기서 \(\overline B\)는 \(B\)가 \(S\)에 대해 갖는 것과 같은 의미를 \(\overline S\)에 대해 갖는다. 따라서 \(R\)은 \(\overline A_r\)에 가약으로 매장되고, \(r=g\bar s\)인 \(\overline A_g\)에 기약으로 매장된다. 여기서는 서로 교환하는 두 계가 각각 \(B,\overline B\)와 동형이다.

\begin{center}
§ 6

\textbf{단순계의 갈루아 이론}
\end{center}

§ 5, 4의 강화된 교환정리는 중심을 기초영역으로 한 단순계의 갈루아 이론을 나타낸다. 먼저 체를, 그 다음 일반 단순계를 고찰한다.

\noindent\textbf{1. 중심에 관하여 유한 랭크인 체의 갈루아 이론.} \(A\)의 \emph{갈루아군} \(\Gcal\)은 중심 \(P\)의 항등사상을 연장하는 모든 자기동형사상의 군, 따라서 § 5, 2에 따라 모든 내부자기동형사상의 군으로 정의한다. 그러므로
\[
\Gcal\simeq A^*/P^*,
\]
여기서 \(A^*,P^*\)는 각각 \(A,P\)의 0 아닌 원소들의 곱셈군이다. 따라서 \(\Gcal\)의 부분군 \(\mathfrak H\)에는 \(P^*\)를 포함하는 \(A^*\)의 부분군 \(H^*\)가
\[
\mathfrak H\simeq H^*/P^*
\]
에 따라 일대일로 대응한다. \(\Gcal\)의 부분군 \(\mathfrak H\)는 \(H^*\)에 0을 덧붙인 것이 한 체 \(H\)가 될 때 닫혀 있다고 한다. \(C\)가 군 \(\mathfrak H\)를 허용한다는 것은 \(C\)와 \(H\)가 원소별로 교환한다는 것과 같다. 따라서 교환정리(§ 5, 4)는 다음 정리가 된다.

\begin{center}
\textbf{\srcspaced{비가환체의 갈루아 이론의 기본정리}}
\end{center}

\emph{\(P\)와 \(A\) 사이의 체 \(C\)들과 \(\Gcal\)의 닫힌 부분군 \(\mathfrak H\)들을 일대일로 대응시켜, \(C\)가 \(\mathfrak H\)의 완전 불변체이고 \(\mathfrak H\)가 \(C\)의 완전 불변군이 되게 할 수 있다.}

\noindent\textbf{\srcspaced{2. 단순계의 갈루아 이론.}}
\(A_f\)가 중심 \(P\)인 단순계이면 갈루아군 \(\Gcal\)은 다시 내부자기동형사상들의 군이고
\[
\Gcal\simeq A_f^*/P^*
\]
이다. 이제 \(A_f^*\)는 \(A_f\)의 정칙원들의 곱셈군을 뜻한다. \(\Gcal\)의 부분군 \(\mathfrak H\simeq H^*/P^*\)는 \(H^*\)가 생성하는 부분환 \(H\)가 단순계이고 \(H^*\)가 \(H\)의 모든 정칙원으로 이루어질 때 단순-닫힌 부분군이라 한다. 교환정리에서 갈루아 이론으로 넘어가는 것은 다음 보조정리로 주어진다.

\noindent\textbf{\srcspaced{보조정리.}} \emph{\(A_f\)에서 \(P\)를 포함하는 각 단순계는 그 정칙원들의 군 \(H^*\)에 의해 생성된다.} \(P\)가 무한히 많은 원소를 가지면 이는 분명하다. 부정원들로 만든 “일반원소”가 정칙이고, 알맞게 특수화하면 \(P\) 위의 정칙 기저원소들을 만들 수 있기 때문이다. Shoda에 따르면 보조정리는 일반적으로도 성립한다\footnote{주 3에서 인용한 Shoda의 논문. 거기서는 부정원의 도입을 피한다.}.

보조정리에 따라 \(S\)와 한 단순계 \(\overline S\)가 교환한다는 것은 다시 \(S\)가 \(\overline S\)의 정칙원들이 유도하는 자기동형사상들을 허용한다는 것과 같다. 따라서 여기서도 교환정리(§ 5, 4)는 다음 정리가 된다.

\noindent\textbf{\srcspaced{단순계의 갈루아 이론의 기본정리.}} \emph{\(A_f\)에서 \(P\)를 포함하는 단순계 \(S\)들과 \(\Gcal\)의 단순-닫힌 부분군 \(\mathfrak H\)들을 일대일로 대응시켜, \(S\)가 \(\mathfrak H\)의 완전 불변영역이고 \(\mathfrak H\)가 \(S\)의 완전 불변군이 되게 할 수 있다.}

\noindent\textbf{\srcspaced{3. 연장원리에 의한 증명.}}
체의 경우에 기본정리의 둘째 증명을 제시한다. 이 증명은 동형, 곧 1차 표현의 연장원리에 바탕을 두며 랭크를 세지 않으므로 더 약한 유한성 가정만 필요하다. 특히 무한 랭크인 체 \(S\)로 어느 방향에서 옮길 수 있는지가 드러난다. 또한 기약 반대 표현류가 하나뿐이라는 사실을 쓰지 않으므로 가환적인 경우에도 그대로 성립한다(§ 8, 2). 먼저 몇 가지 보조정리를 둔다.

\noindent\textbf{\srcspaced{가정.}} \(P\) 위 단순계 \(S\)가 \(A\) 안에 1차 반대 표현을 갖는다고 하자. 따라서 \(S\)는 체이다. 여기서 \(A\)는 중심 \(P\) 위 유한 또는 무한 랭크일 수 있다. \(S_A\)를 \(n\)개의 단순하고 서로 연산자동형인 오른쪽 아이디얼로 분해한 것이(§ 4, 3)
\[
S_A=r_1+\cdots+r_n
=e_1S_A+\cdots+e_nS_A
=e_1A+\cdots+e_nA
\]
로 주어졌다고 하자. \(e_i\)가 생성하는 반대 표현은
\[
e_i s=e_i\sigma
\]
로 정의된다.

\noindent\textbf{\srcspaced{보조정리 1.}} \emph{\(e_1,\ldots,e_n\)이 생성하는 \(n\)개의 반대 표현은 모두 서로 다르다. 따라서 같은 표현류에 속하는 서로 다른 표현들이다. 그러므로 \(S\)는 적어도 그 랭크만큼 서로 다른 표현을 갖는다.}

증명을 위해 § 2의 끝에 따라 \(S\)의 표현들을 \(S_A\)의 대응들로 확장하여 \(A\)에 관하여 연산자준동형이 되게 한다. 여기서는 \(S_A\)의 각 \(w\)에 대하여 \(\omega\in A\)인
\[
e_i w=e_i\omega
\]
로 이 대응을 정의한다. 서로 다른 \(i,j\)에 대하여 \(e_i,e_j\)가 같은 표현을 생성한다면 모든 \(w\in S_A\)에 대해 \(e_iw=e_jw\)이어야 한다. 그러나 \(e_ie_i=e_i\), \(e_je_i=0\)이므로 불가능하다.

\noindent\textbf{\srcspaced{보조정리 2.}} \emph{\(T\)가 \(P,S\) 사이의 체이고 \(S\)의 \(T\)-랭크가 \(s\)이면, \(T\)에서 \(A\)로 가는 각 반대동형은 서로 다른 연장을 적어도 \(s\)개 갖는다.}

\(T\)에서 \(A\)로 가는 반대동형, 곧 \(T\)의 \(A\) 안 반대 표현이
\[
T_A=E_1T_A+\cdots+E_hT_A
=E_1A+\cdots+E_hA,\qquad E_it=E_i\tau
\]
라는 분해로 매개된다고 하자. 여기서 \(E_i\)들은 멱등원이다. 이는 일반성을 제한하지 않는다. 기저원소 \(E_i\alpha\)가 생성하는 표현은 \(\alpha^{-1}E_i\alpha\), 곧 한 멱등원으로도 생성되기 때문이다. 오른쪽에서 \(S\)를 곱하면
\[
S_A=E_1S_A+\cdots+E_hS_A
\]
를 얻는다. 이때 모든 \(E_iS_A\)의 \(A\)-랭크는 \(s\)이다. 실제로 \(S\)에 \(T\)-기저를 대입하고 \(S,A\)의 교환성을 쓰면 랭크가 많아야 \(s\)임을 알 수 있다.
\[
S=Tw_1+\cdots+Tw_s,
\]
\[
E_iS_A=E_iT_Aw_1+\cdots+E_iT_Aw_s
=E_iw_1A+\cdots+E_iw_sA.
\]
합 \(S_A\)의 랭크가 \(n=hs\)이므로 각 \(E_iS_A\)의 랭크는 정확히 \(s\)이다. 따라서
\[
E_iS_A=r_{i1}+\cdots+r_{is}
=e_{i1}A+\cdots+e_{is}A,
\]
이고 보조정리 1에 따라 \(e_{i1},\ldots,e_{is}\)가 생성하는 \(s\)개의 표현은 모두 서로 다르다. 이들은 모두 \(E_i\)가 생성하는 \(T\)의 표현의 연장이다. 실제로 \(E_it=E_i\tau\)에서
\[
(e_{i1}+\cdots+e_{is})t=(e_{i1}+\cdots+e_{is})\tau
\]
가 나오고, 직합분해 때문에 \(e_{ij}t=e_{ij}\tau\)이다.

\paragraph{정의.}
\(T\)가 유도하는, \(S\)에서 \(A\)로 가는 반대동형들의 계 \(\mathfrak I\)의 류분할 \(\mathfrak I_T\)를 다음처럼 정의한다. \(T\) 위에서 같은 동형을 유도하는 \(\mathfrak I\)의 동형들과 오직 그것들만 서로 동치라고 한다.

\noindent\textbf{\srcspaced{기본정리.}} \emph{\(\mathfrak I_T\)가 \(T\)가 유도한 류분할이면, \(T\)는 이 류분할에 관한 극대 부분체이다.}

실제로 \(L\)이 \(T,S\) 사이의 중간체이고 \(T\) 위 차수가 \(l\)이면, 보조정리 2에 따라 \(\mathfrak I_T\)의 각 류는 \(\mathfrak I_L\)의 류 적어도 \(l\)개로 갈라진다. 이 정식에서 통상적인 기본정리로 넘어가려면, 동형들의 집합 \(\mathfrak I\)의 원소들을 그 가운데 고정한 하나의 역원으로 곱하여 자기동형사상군으로 옮기면 된다. 이때 \(\mathfrak I_T\)의 한 류는 \(T\)의 불변군으로 옮겨간다. 닫힌 부분군들이 완전 불변군이라는 명제도 여기에 포함되어 있다. \(\mathfrak H\simeq H^*/P^*\)에 대해서는 체 \(H\)를 중간체 \(T\)로 대입하기만 하면 된다.

\begin{center}
§ 7

\textbf{\srcspaced{유사한 대수들의 류. 분해체}}
\end{center}

이제부터는 중심 \(P\)에 관하여 \emph{유한 랭크}인 체 \(A,\overline A\)와 행렬환 \(A_r\)만 다룬다. 곧 \(P\) 위의 단순 정규대수이며, 이를 줄여서 \emph{\(P\) 위의 대수} 또는 단순히 대수라 하겠다. \(A,\overline A\)는 이때 \emph{나눗셈대수}이다. 한 \emph{유사한 대수들의 류}, \(A_r\sim A_s\), 에는 같은 딸린 \(A\)를 갖는 모든 \(A_r\)를 묶는다. 동형인 \(A\)들은 동일시한다.

\noindent\textbf{\srcspaced{1. 대수류들의 군.}} \(A_r,B_s\)가 \(P\) 위의 대수이면 그 곱, 곧 \(P\) 위의 직접곱도 그러하다. § 3의 끝에 따라
\[
A_r\times B_s=C_t
\]
이고, 여기서 \(C\)는 동형까지 유일하게 정해지는 중심 \(P\)인 나눗셈대수이다. \(A_r,B_s\)에 \(P\) 위 행렬환들을 곱해 보면 이 곱셈이 류들에 대해 일의적으로 정해짐을 알 수 있다. 따라서 류들은 곱셈에 닫힌 가환계를 이룬다. 더 나아가 다음이 성립한다.

\noindent\textbf{\srcspaced{R. 브라우어의 정리.}} \emph{유사한 대수들의 류는 직접곱에 관하여 아벨군을 이룬다.}

실제로 이 계는 \(P\) 위의 행렬환들, 곧 \(P\)와 유사한 대수들로 이루어진 항등류를 갖는다. 또한 각 류 \((A)\)에 역류 \((A)^{-1}\)가 있으며, 이는 \(A\)와 반대동형인 \(\overline A\)에 유사한 대수들로 이루어진다. 이는 § 5, 3의 교환정리에서 \(S=A\)로 특수화하면 나온다. 실제로 \(A\)는 자기 자신 안에 기약으로 매장되므로 \(B=P\), \(A_A=P_t\)이다\footnote{대수류군에 관한 더 세밀한 정리들은 인자계 이론을 쓴다. 여기서는 다루지 않는다. 여기까지는 가환체에 관한 고찰이 전혀 나타나지 않았음을 지적한다. 예컨대 랭크 \((A:P)\)가 제곱수라는 사실은 사용하지도 증명하지도 않았다.}.

\noindent\textbf{\srcspaced{2. 한 대수류의 분해체.}}\quad 분해체 이론도 전적으로 § 5의 처음에 말한 원리에 바탕을 둔다. 가환 확대체들을 \(A\) 또는 \(\overline A\) 안에 표현한다. 먼저 다음 보조정리를 둔다.

\medskip
\noindent\textbf{\srcspaced{보조정리.}} \emph{\(A\)가 중심 \(P\)인 체이고 \(Z\)가 \(P\)의 유한 가환 확대체이면 \(A_Z\)는 중심 \(Z\)인 단순계이다.} 실제로 § 4, 1에 따라 \(A_Z\)와 반대동형인 계 \(Z_A\)에 이 사실이 성립한다. 따라서 \(P\) 위에서 \(A\)와 유사한 대수들의 각 류 \((A)\)에 대하여, \(Z\)는 \(Z\) 위에서 \(A_Z\)와 유사한 단순대수들의 \emph{확대류} \((A)_Z\)를 생성한다.

확대류에 딸린 나눗셈대수 \(D\)는 § 5, 3의 교환정리에서 곧바로 나온다.

\noindent\textbf{\srcspaced{확대류에 딸린 나눗셈대수에 관한 정리.}} \emph{\((A)_Z\)가 확대류이고 \(Z\)가 \(A_r\) 안에 기약으로 매장, 곧 표현되었다고 하자. 그러면 \((A)_Z=(D)\)이며, \(D\)는 \(A_r\)에서 \(Z\)의 모든 원소와 교환하는 원소들의 전체이다.} § 5, 3의 단순계 \(S\)를 \(Z\)로 바꾸고 \(Z_A,A_Z\)가 반대동형임을 고려하면 된다. \(Z\)가 가약으로 매장되었으면 \(Z\)의 모든 원소와 교환하는 원소 전체는 \(D\) 위의 행렬환이 된다.

\(P\)의 가환 확대체 \(Z\)는 확대류 \((A)_Z\)가 완전히 분할될 때, 곧 \(Z\) 위의 항등류, 즉 \(Z\) 위 모든 행렬환의 류와 같을 때 류 \((A)\)의 \emph{분해체}라 한다.

\noindent\textbf{\srcspaced{분해체의 특성화 정리.}}
\emph{\(P\)의 유한 가환 확대체 \(Z\)가 류 \((A)\)의 분해체인 것은, \(Z\)를 \(A_r\) 안에 기약으로 매장한 것이 \(A_r\)의 극대 가환 부분체가 될 때 그리고 오직 그때이다.}

\(Z\)가 \((A)\)의 분해체라는 성질은 \((D)=(Z)\)와 같다. 따라서 딸린 나눗셈대수에 관한 정리에 따라 기약으로 매장된 \(Z\)는 극대 가환 부분체여야 한다. 역으로 이 조건이 성립하면 반드시 \(D=Z\)이다. 실제로 \(D\)의 각 원소 \(a\)를 \(Z\)에 첨가하면 가환 확대체를 생성한다.

\noindent\textbf{\srcspaced{주의.}}\quad
1. 동시에 \emph{극대 가환 부분체} \(Z\)를 \emph{기약으로} 매장하면 \emph{극대 가환 부분환}을 생성함이 따라온다. \(Z\)의 모든 원소와 교환하는 원소 전체가 체를 이루기 때문이다.

2. 이 정리는 분해체의 \emph{존재}도 준다. 딸린 나눗셈대수 \(A\)는 반드시 극대 가환 부분체를 갖기 때문이다.

\medskip
\noindent\textbf{\srcspaced{3. 랭크 관계, 인덱스, 무한 가환 확대.}}\quad § 5, 4의 랭크 명제는 먼저 단순대수의 중심에 관한 랭크가 제곱수라는 알려진 사실을 준다. \(Z\)가 \(A\)의 극대 가환 부분체이면 \(A\) 안의 모든 매장은 기약이므로
\[
(Z:P)(Z:P)=(A:P).
\]
따라서
\[
(A:P)=m^2,\qquad(Z:P)=m,
\]
이고 \(m\)은 \emph{슈어 인덱스}이다. 이는 동시에 나눗셈대수 \(A\) 자체에 매장 가능한 모든 분해체의 차수이다. 일반 분해체의 차수 \(n\)에 대해서는
\[
n=m\cdot r
\]
을 얻는다\footnote{일반적으로는 진부분체가 분해체가 아닌 “최소” 분해체가 모든 차수에 걸쳐 존재하기도 한다. 주 2에서 인용한 Brauer--Noether 논문을 보라.}. 예를 들어 \((A_r:P)=m^2r^2\)에서 나오며, § 4, 3의 랭크 관계에서도 나온다. 후자의 이유는 인덱스 \(m\)이 \(A_Z\)의 단순성분 수와도 같고, \(A_Z\)가 \(Z\) 위 랭크 \(m^2\)인 행렬환이 되기 때문이다. 더 일반적으로 \(A_\Lambda\sim D\)이고 \(L\)이 \(\Lambda\)의 \(A_r\) 안 기약 매장이며 나눗셈대수 \(D\)의 중심 \(L\) 위 랭크 \((D:L)\)가 \(d^2\)이면
\[
l\cdot d=m\cdot r.
\]
실제로 § 5, 4와 딸린 나눗셈대수에 관한 정리에 따라
\[
(A_r:P)=(L:P)(D:P),
\quad\text{즉}\quad m^2r^2=l^2d^2
\]
이다.

분해체의 존재로부터 \(P\)의 무한 가환 확대체 \(\Omega\), 대수적이든 초월적이든, 에 대하여 다음이 따라온다.

\noindent\emph{확대류 \((A)_\Omega\)는 중심 \(\Omega\)인 단순대수들의 류이다. 특히 \(\Omega\)가 대수적으로 닫혀 있으면 \(A_\Omega\)는 \(m\)차 행렬환이다. 따라서 인덱스는 “절대 성분수”와 같다.}

실제로 \(Z\)가 \(A\)의 분해체이고 \(\bar\Omega\)가 \(\Omega,Z\)의 합성체이면 \(A_{\bar\Omega}\)는 \(\bar\Omega\) 위 행렬환이다. 따라서 라디칼이 없고 중심이 \(\bar\Omega\)이다. 그러므로 \(A_\Omega\)도 라디칼이 없으며 중심은 \(\Omega\)이다. 실제로 \(A_\Omega\)의 \(\Omega\) 밖 원소가 중심에 있으면 \(A_{\bar\Omega}\)의 중심에도 있게 된다. 따라서 \(A_\Omega\)는 \(\Omega\) 위 단순대수이고, \(\Omega\)가 대수적으로 닫혀 있으면 반드시 \(\Omega\) 위 행렬환이다.

\medskip
\noindent\textbf{\srcspaced{4. 분리가능 분해체의 존재.}}\quad
\emph{모든 류 \((A)\)는 분리가능 분해체를 가지며, 심지어 \(A\) 자체에 매장 가능한 것을 갖는다.}\footnote{G. Köthe, \emph{Über Schiefkörper mit Unterkörpern zweiter Art über dem Zentrum}, Journ. f. Math. 166 (1932), 182--184쪽을 보라. 여기의 훨씬 간단한 증명은 M. Zorn의 지적에서 나왔다.}

기초체 \(P\)의 표수를 \(p\), 인덱스를 \(m=s\cdot p^f\)라 하자. 여기서 \(s\)는 \(p\)와 서로소이다. \(Z\)가 \(A\)의 차수 \(m\)인 분해체이고 \(Z_0\)가 그 안에 든 제1종 확대, 따라서 \((Z:Z_0)=p^f\), 이면 3에 따라 \(A_{Z_0}\sim D\)의 인덱스는 \(p^f\)이다. \(Z_0\)의 분리가능 확대체 \(Z_1\) 가운데 \(D_{Z_1}\)의 인덱스가 \(p^f\)의 진약수가 되는 것이 존재함을 보이겠다. 대수적으로 닫힌 \(\Omega\)에 대한 \(D_\Omega\)는 3에 따라 라디칼이 없으므로 \(D\)의 축약판별식은 0이 아니다. 따라서 \(D\)에는 축약대각합이 0이 아닌 원소 \(d\)가 적어도 하나 있다. 이 \(d\)는 기초체 \(Z_0\)에 속할 수 없다. \(Z_0\)의 각 원소 \(\alpha\)의 대각합은 \(p^f\alpha=0\)이기 때문이다. 그러므로 \(Z_1=Z_0(d)\)는 진확대이고 분리가능 확대체이며, \(D_{Z_1}\)의 인덱스는 \(p^f\)의 진약수이다. 이를 유한 번 반복하면 인덱스와 같은 차수 \(m\)을 갖는 \((A)\)의 분리가능 분해체에 이른다.

\begin{center}
\textbf{§ 8}\\[0.35em]
\textbf{가환적인 경우의 분해체·분리체·갈루아 이론}
\end{center}

중심 위 단순계의 분해체에서 임의의 단순계의 분해체로 넘어가려면 중심의 분해이론을 더 전개해야 한다. 여기에 § 6, 3과 평행한 갈루아 이론도 덧붙인다. 이 절의 모든 체와 계는 가환이다.

\noindent\textbf{\srcspaced{1. 가환 단순계의 분해체와 분리체.}}\quad 가환계 \(Z\)가 \(P\) 위에서 단순, 따라서 체이고 \(\Omega\)가 \(P\)의 대수적으로 닫힌 확대체라 하자. 그러면 다음은 알려져 있다(\emph{표현론} § 21).

\noindent\emph{\(Z_\Omega\)가 라디칼 없는 계로 남는 것은 \(Z\)가 \(P\) 위 분리가능, 곧 제1종 확대일 때 그리고 오직 그때이다.}

실제로 이때 그리고 오직 이때 \(Z_\Omega\)는 \(\Omega\) 위 랭크와 같은 수의 1차 동형사상을 \(\Omega\) 안에 가지며, 따라서 같은 수의 서로 다른 표현을 갖는다. 여기서는 표현과 표현류가 일치한다.

\(Z_\Omega\)에 라디칼이 나타날 수 있고 따라서 절대단순 성분들로의 직합분해가 반드시 존재하지는 않는다는 사실은 다음 정의들을 낳는다.

\noindent\textbf{\srcspaced{정의.}}\quad \(P\)의 확대체 \(\Lambda\)는 \(Z_\Lambda\)의 합성인자들이 모두 1차일 때, 달리 말해 모든 절대기약 표현이 이미 \(\Lambda\) 안에 있을 때 분해체라 한다. \(P\)의 확대체 \(T\)는 \(Z_T\)가 적어도 하나의 1차 합성인자를 분리해 낼 때, 달리 말해 \(Z\)의 절대기약 표현이 적어도 하나 \(T\) 안에 존재할 때 분리체라 한다.

분해체와 분리체는 다음 정리로 특성화된다.

\noindent\textbf{\srcspaced{정리.}} \emph{한 체 \(T\)가 분리체인 것은 \(Z\)와 동형인 부분체를 포함할 때 그리고 오직 그때이다. 한 체 \(\Lambda\)가 분해체인 것은 \(Z\)에 딸린 갈루아체와 동형인 부분체 \(\Gamma\)를 포함할 때 그리고 오직 그때이다.}

이는 절대기약 표현을 이용한 분해체와 분리체의 정의에서 바로 나온다. 특히 비가환적인 경우와 유사하게 다음 귀결을 얻는다.

\noindent\textbf{\srcspaced{귀결.}} \emph{한 체가 \(Z\)와 동형일 때 그리고 오직 그때 \(Z\)의 최소 분리체, 곧 진부분체가 분리체가 아닌 체이다. 이는 그 체를 \(Z\)에 기약 가환적으로 매장한 것이 \(Z\)의 극대 가환 부분체를 이룬다는 뜻이다. 최소 분리체가 동시에 분해체인 것은 \(Z\)가 \(P\) 위 정규체, 곧 갈루아체일 때 그리고 오직 그때이다.}

단순대수를 그 중심 위에서 정규라 부르는 이유가 여기에 있다. 고정한 \(P\) 위 대수적으로 닫힌 \(\Omega\) 안에서는 최소 분해체가 \(Z\)에 딸린 갈루아체 하나만 존재한다. 이는 일반적으로 동형이 아닌 최소 분해체가 무한히 많은 비가환적인 경우와 대조된다\footnote{R. Brauer--E. Noether, 앞의 글.}. 그 이유는 가환적인 경우 직합분해가 유일한 반면 비가환적인 경우에는 연산자동형까지 유일할 뿐이라는 데 있다. 같은 말로 하면, 가환적인 경우에는 서로 다른 절대기약 표현이 유한개뿐인 반면 비가환적인 경우에는 같은 류에 속하면서도 서로 다른 것이 무한히 많기 때문이다(2를 보라).

\medskip
\noindent\textbf{\srcspaced{2. \(Z/P\)가 분리가능체일 때 갈루아 이론의 초복소적 정초.}}\quad § 6, 3과 정확히 유사하게, \(P\) 위 분리가능인 임의의 \(Z\)에 대해 동형 이론이 성립한다. \(Z\)가 갈루아체인 경우에는 통상적인 자기동형사상군으로 옮길 수 있다.

\medskip
\noindent\textbf{\srcspaced{가정.}}\quad \(P\) 위 단순계 \(Z\)를 유한 분리가능 확대체라 하고, \(\Omega\)를 \(P\) 위 유한 또는 무한 분해체라 하자. \(Z=Z^{(1)}\)은 \(Z\)와 동형인 부분체, \(\Gamma\)는 그에 딸린 갈루아체라 하자. 1에 따라 \(Z\)가 분리가능이므로 직합분해
\[
Z=r^{(1)}+\cdots+r^{(n)}
=e^{(1)}Z_\Omega+\cdots+e^{(n)}Z_\Omega
=e^{(1)}\Omega+\cdots+e^{(n)}\Omega
\]
가 성립한다. \(e^{(i)}\)가 생성하는 표현 \(Z\to Z^{(i)}\), \(Z^{(i)}\leqq\Gamma\), 는
\[
e^{(i)}z=e^{(i)}\zeta^{(i)}
\]
로 정의된다.

\noindent\textbf{\srcspaced{보조정리 1.}} \emph{\(e^{(1)},\ldots,e^{(n)}\)이 생성하는 \(n\)개 표현은 모두 서로 다르다.} 증명은 § 6, 3과 같고, 1에서도 나온다. 이들은 서로 다른 표현류에 속하기 때문이다.

\noindent\textbf{\srcspaced{보조정리 2.}} \emph{\(T\)가 \(P,Z\) 사이의 중간체이고 \(Z\)의 \(T\)-랭크가 \(s\)이면, \(T\)에서 \(\Omega\)로 가는 각 동형은 적어도 \(s\)개, 따라서 정확히 \(s\)개의 연장을 갖는다.}

증명은 § 6, 3과 같다. 여기서 “적어도”를 “정확히”로 강화할 수 있는 것은 \(Z_\Omega\)의 분해가 유일하기 때문이다. 이에 따라 \(Z\)는 \(\Omega\) 안에 적어도 \(n\)개가 아니라 정확히 \(n\)개의 동형을 갖는다(1 끝의 주의를 보라).

§ 6, 3과 마찬가지로 이제 \(Z\)의 동형들로 이루어진 유한집합 \(\mathfrak J\)에 \(T\)가 \emph{유도하는 류분할} \(\mathfrak J_T\)를 정의한다. 곧 \(T\) 위에서 같은 동형을 유도하는 \(\mathfrak J\)의 동형들과 오직 그것들만 \(\mathfrak J_T\)에서 서로 동치라고 한다. 그리고 다음을 증명한다.

\noindent\textbf{\srcspaced{기본정리.}} \emph{\(\mathfrak J_T\)가 \(T\)가 유도하는 류분할이면 \(T\)는 이 류분할에 관한 극대 부분체이다.}

\(Z\)가 갈루아체인 경우에는 § 6, 3과 똑같이 한 동형의 역과 합성하여 자기동형사상군 및 통상적인 정식으로 옮길 수 있다. 부분군에 관한 대응 명제는 그 자체로 흥미로운 멱등원을 고찰하면 나온다.

\medskip
\noindent\textbf{\srcspaced{3. \(Z_\Omega\)의 멱등원.}}\quad \(S\)가 \(Z\)를 각각의 \(Z^{(i)}\)로 옮기는 \(n\)개 치환을 두루 돌게 하자. 이는 \(Z\to Z\)와 \(Z\to Z^{(i)}\)라는 동형들의 합성이며, 첫 것은 \(Z\to Z\)의 역동형을 뜻한다. \(Z\)는 갈루아체일 필요가 없다. \(Z_Z\)의 원소 \(a\)에 대하여 치환 \(S\)를 연산자로 정의하되, \(Z\) 위에서는 항등사상을 유도한다고 정한다. 그러면 각 \(a\)에 대해 \(Z_{Z^{(i)}}\)에 놓이는 켤레원소 \(a^S\)가 정의된다. 이 규정 아래 다음이 성립한다.

\noindent\textbf{\srcspaced{정리.}} \emph{\(Z_\Omega\)의 \(n\)개 멱등원 \(e^{(i)}\)는 서로 켤레이다. \(e=e^{(1)}\)라 하면 \(e^{(i)}=e^S\)이다. 이들은 \(n\)개의 켤레 확대계 \(Z_{Z^{(i)}}\)에 각각 놓인다.}

실제로 \(S\)를 적용하면 정의관계 \(ez=e\zeta^{(1)}\), \(e^2=e\)는
\[
e^Sz=e^S\zeta^{(i)},\qquad(e^S)^2=e^S
\]
로 옮겨간다. 분해의 유일성 때문에 멱등원들도 유일하게 정해지므로 \(e^{(i)}=e^S\)이다. 또한 \(Z\)는 분리체여서 표현 \(ez=e\zeta\)를 이미 매개하므로 \(e\in Z_Z\), 따라서 \(e^{(i)}\in Z_{Z^{(i)}}\)이다.

특히 \(Z\)가 갈루아체이면 \emph{갈루아 이론 기본정리에서 부분군에 관한 부분}이 곧바로 따라온다. 실제로 \(\mathfrak H\)가 갈루아군 \(\Gcal\)의 부분군이면, 직합성분이 유일하게 정해진다는 데서 \(H\in\mathfrak H\)에 대한 원소 \(\sum e^H\)는 \(\mathfrak H\)의 모든 치환과 오직 그것들만 허용한다. 군 \(\Gcal\)은 \(Z\)에 대해 정의했으므로 이는 \(Z\)의 갈루아 이론이며, 동형에 의해 \(Z\)에 관한 대응 결과도 함께 얻는다.

\paragraph{4. 멱등원과 여차기저의 관계.}
다시 \(Z\)가 반드시 갈루아체라고 가정하지 않는다.

\noindent\textbf{\srcspaced{정리.}} \emph{\(a_1,\ldots,a_n\)이 \(Z/P\)의 기저이고 멱등원을 \(e=a_1\beta_1+\cdots+a_n\beta_n\), 따라서 \(e^S=a_1\beta_1^S+\cdots+a_n\beta_n^S\), 로 두면, \(\beta_1^S,\ldots,\beta_n^S\)는 \(e^S\)가 매개하는 사상에서 \(a_1,\ldots,a_n\)의 여차기저 \(b_1,\ldots,b_n\)의 상이다.}

실제로 \(e^S\)가 매개하는 사상
\[
e^Sz=e^S\zeta^S,
\]
특히 \(e^Sa_i=e^S\alpha_i^S\)에서 \(e^Se^S=e^S\), \(e^Se^R=0\) (\(S\ne R\))이므로 대응 \(e^S\mapsto1\), \(e^R\mapsto0\)이 성립한다. 따라서
\[
1=\alpha_1^S\beta_1^S+\cdots+\alpha_n^S\beta_n^S,\qquad
0=\alpha_1^S\beta_1^R+\cdots+\alpha_n^S\beta_n^R\quad(S\ne R),
\]
곧 행렬로 쓰면 여차기저의 정의방정식
\[
\begin{pmatrix}
\alpha_1&\cdots&\alpha_n\\
\cdots&\cdots&\cdots\\
\alpha_1^S&\cdots&\alpha_n^S\\
\cdots&\cdots&\cdots
\end{pmatrix}
\begin{pmatrix}
\beta_1&\cdots&\beta_1^R&\cdots\\
\cdots&\cdots&\cdots&\cdots\\
\beta_n&\cdots&\beta_n^R&\cdots
\end{pmatrix}
=E
\]
을 얻는다. 대각합을 통한 정의로 옮기는 것은 다음에서 나온다. \(a_i=\sum_S e^S\alpha_i^S\), \(b_i=\sum_S e^S\beta_i^S\)에 대하여 행렬관계의 두 행렬을 바꾸면
\[
\Sp(a_i b_i)=1,\qquad
\Sp(a_i b_k)=0\quad(i\ne k)\srcfnmark{21)}
\]
가 된다.
\srcfntext{21)}{\emph{표현론} § 25를 보라.}
이때 \(a_i,b_i\)는 모든 치환 \(S\)를 허용하므로 \(Z\)에 놓인다(§ 3, 3의 보조정리)\srcfn{22)}{여차기저와 멱등원 사이의 관계는 Dedekind, \emph{Zur Theorie der aus \(n\) Haupteinheiten gebildeten komplexen Größen}에 처음 나타난다. \emph{Gesammelte Werke} II, 4--5쪽을 보라.}.

여차기저들이 반대변환한다는 사실은 멱등원 \(e^S\)가 \(Z\)의 불변량이고 바탕으로 삼은 기저와 무관하다는 데서 나온다. 곧 \(Z/P\)의 다른 기저 \(\bar a_1,\ldots,\bar a_n\)으로 옮겨갈 때 모든 \(a_1\beta_1^S+\cdots+a_n\beta_n^S\)는 자기 자신으로 옮겨간다.

\begin{center}
\textbf{§ 9}\\[0.8ex]
\textbf{임의의 계의 분해체와 분리체}
\end{center}

\noindent\textbf{\srcspaced{1. 정의와 단순한 경우로의 환원.}}
§ 8의 정의들에 대응하여 일반적인 경우에는 다음과 같이 정의한다.

\noindent\textbf{\srcspaced{정의.}} \(S\)가 \(P\) 위 초복소계라 하자. \(P\)의 가환 확대체 \(\Lambda\)는 \(S_\Lambda\)가 한쪽 아이디얼의 합성열에서 절대단순 합성인자들로 분해될 때, 달리 말해 \(S\)의 \(\Lambda\) 안 모든 기약 표현이 이미 절대기약일 때 \(S\)의 분해체라 한다. \(P\)의 확대체 \(T\)는 \(S_T\)가 적어도 하나의 절대단순 합성인자를 분리해 낼 때, 곧 \(T\) 안의 기약 표현 적어도 하나가 이미 절대기약일 때 분리체라 한다.

이 정의들은 § 8에서 가환적인 경우에, § 7에서 단순대수에 대해 준 정의들을 분명히 포함한다. 후자에 대해서는 중심을 가환적으로 확대할 때마다 \(A_\Lambda\)가 완전가약이므로 합성인자와 직합분해의 성분이 일치하기 때문이다. 여기서 중심 위 완전행렬환과 오직 그것들만 절대단순 성분, 따라서 절대기약 표현을 준다. 또 \(A_\Lambda\)는 양쪽 단순이어서 서로 연산자동형인 성분들로 분해되므로 절대단순 성분 하나를 분리해 내기만 해도 완전히 분할된다. 따라서 분리체와 분해체가 일치한다.

정의에서 다음 사실들이 바로 나온다. \emph{\(S\)의 분해체는 \(P\)에서 기약인 각 표현의 분해체 하나씩을 합성한 체들이다. \(S\)의 분리체는 \(P\)에서 기약인 어떤 표현의 분리체이면 된다.} 단순계, 곧 라디칼로 나눈 잉여류환의 성분들과 \(P\) 안의 기약 표현들이 일대일로 대응하므로 \emph{단순계만 고찰하면} 충분하다. 이 경우 결과들은 § 7과 § 8을 결합하여 얻는다.

\noindent\textbf{\srcspaced{2. 단순계의 분해체와 분리체.}} 먼저 중심 \(Z\)가 \(P\) 위 \emph{분리가능}인 단순계 \(A\)를 고찰한다. § 8, 1과 § 8, 3에서 다음이 나온다.

\emph{\(\Omega\)가 \(Z\)의 분해체일 때 \(Z_\Omega\)의 직합분해에 대응하는 \(A_\Omega\)의 양쪽 분해는 \(A\)와 동형인 \(n\)개의 켤레 성분 \(e^{(i)}A\)를 준다. 이 성분들은 각자의 중심 \(e^{(i)}Z^{(i)}=e^{(i)}Z\) 위 단순대수가 된다. 여기서 치환 \(S\)는 \(A\) 위 항등사상을 유도하도록 규정하여 \(A_Z\)의 연산자로 정의한다.} (§ 8, 3을 보라.)

실제로 § 8, 3에 따라 멱등원들이 켤레이므로, \(A\)에 대한 치환의 정의상 성분 \(e^{(i)}A\)들도 켤레이다. \(A\)가 양쪽 단순이므로 \(e^{(i)}A\)는 \(A\)와 환 동형이다. 또한 \(e^{(i)}Z=e^{(i)}Z^{(i)}=e^{(i)}Z_{Z^{(i)}}\)이므로
\[
e^{(i)}A=e^{(i)}A_{Z^{(i)}}.
\]
따라서 중심은 계수영역과 일치하며, 성분들은 정규대수이다.

이제 § 7, 2의 결과들은 다시 다음을 준다.

\emph{\(A\)가 \(P\) 위 단순계이고 중심 \(Z\)가 \(P\) 위 분리가능이면 \(A_\Omega\)도 라디칼이 없고 따라서 완전가약이다. 여기서 \(\Omega\)는 임의의 유한 또는 무한 확대체일 수 있다.}

실제로 § 7, 2에 따라 \(e^{(i)}A_{Z^{(i)}}\)는 계수를 어떻게 확대해도 단순이고 따라서 라디칼이 없다. 그러므로 직합도 라디칼이 없다. 이 직합은 \(A_\Omega\)이거나, \(\Omega\)가 \(Z\)의 분해체가 아니면 \(A_\Omega\)에서 계수를 확대하여 생긴다. § 7, 2로부터 다음도 나온다.

\begin{center}
\textbf{\srcspaced{분리체와 분해체의 특성화}}
\end{center}

\emph{\(A\)가 분리가능 중심 \(Z\)를 갖는 나눗셈대수이면, 분리체의 기약 매장은 \(A_r\)에서 \(Z\)를 포함하는 극대 가환 부분체들과 정확히 일치한다. 분해체는 분리체 하나와 중심의 분해체 하나, 특히 중심에 딸린 갈루아체를 합성한 체들과 정확히 일치한다. 중심이 갈루아체가 아니어도 최소 분리체가 분해체일 수 있다.}

실제로 분리체는 \(Z\)와 동형인 체를 포함해야 한다. 분리체의 매장에 관한 명제는 이제 앞의 사실들과 § 7, 2에서 나온다. 분해체는 더 나아가 \(Z\)의 한 분해체 \(\Omega\)를 포함해야 한다. 그런데 \(\Omega\) 위의 모든 분리체는 동시에 분해체이다. 성분 \(e^{(i)}A_\Omega\)가 단순이고 \(\Omega\)와 동형인 중심을 갖기 때문이다. 따라서 \(Z\)가 갈루아체이면 최소 분리체와 최소 분해체가 일치한다. 그러나 가환적인 경우와 달리 비갈루아인 경우에도 최소 분리체가 동시에 분해체일 수 있다. 바로 그러한 최소 분리체가 \(Z\)에 딸린 갈루아체를 포함하는 경우이다. 예를 들어 \(P\)가 유리수체일 때 \(P(\sqrt[3]{2})\)와 동형인 중심을 갖는 사원수체에서는 \(P(\sqrt[3]{2})\)에 딸린 갈루아체가 최소 분리체이자 분해체이다.

끝으로 중심이 \emph{비분리가능이면 \(Z_\Omega\), 따라서 \(A_\Omega\)는 라디칼을 갖는 계이다.} 예를 들어 \(\mathfrak C\)를 \(Z_\Omega\)의 라디칼, \(\mathfrak c\)를 \(A_\Omega\)에서 그 확대아이디얼이라 하자. 그러면 \(Z_\Omega/\mathfrak C\)에는 § 8, 2와 정확히 유사하게 켤레 성분들로의 직합분해가 성립한다. 위와 똑같이 \(A_\Omega/\mathfrak c\)의 직합분해를 얻으며, 그 성분 \(\bar e^{(i)}A\)들은 \(A\)와 동형이고 중심은 \(\bar e^{(i)}Z^{(i)}\)이다. 따라서 분리체와 분해체에 관한 정리들도 유지된다. 더 나아가 랭크를 세면 \(\mathfrak C\)의 합성열에 대응하는 합성인자들이 \(A_\Omega/\mathfrak c\)와 단순히 준동형일 뿐 아니라 연산자동형임을 알 수 있다.

기약 표현의 말로 쓰면 다음과 같이 요약된다.

\noindent\textbf{\srcspaced{요약.}} \emph{한 초복소계의 \(P\) 안 기약 표현이 주어졌을 때, 이 표현이 절대 완전가약인 것은 그에 딸린 중심이 \(P\) 위 분리가능일 때 그리고 오직 그때이다. 어느 경우든 표현의 분해체와 분리체는 위와 같이 딸린 단순계 안의 매장으로 특성화할 수 있다. 특히 분리체의 최소차수는 중심의 차수와 중심 위 대응 대수류의 인덱스의 곱과 같고, 모든 분리체의 차수는 그 배수이다. 이 표현은 합성열의 뜻에서, 중심에 포함된 가장 큰 분리가능체의 차수만큼 서로 다르지만 서로 켤레인 절대기약 표현류들로 분해된다. 각 류는 인덱스와 그 분리가능 확대에 대한 중심의 차수의 곱만큼 반복된다.}

모든 확대체가 분리가능인 완전체의 경우에는 이렇게 I. Schur의 알려진 결과들을 다시 얻으며, R. Brauer에게 이미 있는 매장에 의한 특성화를 덧붙이게 된다.

\begin{center}
(1932년 6월 8일 접수)
\end{center}
